% Figure: a real-ish tensile test dataset (data/aluminum_6061_tensile.csv) with
% the elastic-region linear fit giving Young's modulus. Plots stress vs strain
% in the low-strain region and the fitted line.
\begin{figure}[htbp]
\centering
\begin{tikzpicture}
\begin{axis}[
  width=12cm, height=7.5cm,
  xlabel={engineering strain $\epsilon$ ($\times 10^{-3}$)},
  ylabel={engineering stress $\sigma$ (MPa)},
  xmin=0, xmax=6, ymin=0, ymax=320,
  axis lines=left,
  tick label style={font=\scriptsize},
  label style={font=\small},
  legend pos=south east,
  legend style={font=\scriptsize},
]
% measured points (elastic region of a 6061-T6 coupon), strain in 1e-3 units
\addplot[only marks, engblue, mark=*, mark size=1.4pt] coordinates {
  (0.00,0) (0.50,34) (1.00,69) (1.50,104) (2.00,138)
  (2.50,172) (3.00,205) (3.50,238) (3.90,262) (4.40,278)
  (5.00,288) (5.60,295)
};
\addlegendentry{measured (Al 6061-T6)}
% linear fit through elastic region: slope ~ 69 GPa => 69 MPa per 1e-3 strain
\addplot[engred, very thick, samples=2, domain=0:3.7] {69*x};
\addlegendentry{elastic fit, $E\approx 69$ GPa}
% mark yield (0.2% offset, schematic)
\filldraw[engorange] (axis cs:4.4,278) circle [radius=2pt];
\node[engorange, font=\scriptsize, anchor=north west] at (axis cs:4.4,278) {0.2\% offset yield $\approx 276$ MPa};
\end{axis}
\end{tikzpicture}
\caption[Tensile-test data and the modulus fit]{Engineering stress
versus strain for a 6061-T6 aluminium coupon (data in
\texttt{data/aluminum\_6061\_tensile.csv}). The slope of the
straight-line elastic region is Young's modulus; a least-squares fit
through the first six points returns $E\approx 69$ GPa, the textbook
value for this alloy. The departure from the line above about
$\epsilon=0.004$ marks the onset of yield; the 0.2 percent offset
construction (a line of slope $E$ shifted to $\epsilon=0.002$) gives
a yield stress of roughly $276$ MPa. The fit is computed in
\texttt{code/stress\_strain\_fit.py}.}
\label{fig:vol03ch03:tensile-test-fit}
\end{figure}
